% ---------------------------------------------------------------------------
%  III. Platform and Access Model
%  Floats:  figures/fig-network.tex   (Fig. 2, three network segments)
%  Policy:  the robot model is deliberately not named. Use \platform{} /
%           \Platform{} from preamble.tex, never a product reference.
% ---------------------------------------------------------------------------

\section{Platform and Access Model}
\label{sec:access}

\Platform{} is a mass-produced indoor reception robot, sold as a turnkey product with a
companion Android application and a subscription cloud service. It pairs a Linux/\ros{}
chassis~\cite{quigley2009ros} carrying SLAM, a planner, LiDAR and RGB-D sensing with an
Android~7.1 head that drives a \SI{15.6}{inch} touchscreen, cameras and speakers. We refer to it
generically throughout, since nothing in the methodology depends on the particular model and we
have no interest in singling out one manufacturer. The two computers run independent network
stacks, a fact no vendor document states and no application string reveals.

% Fig. 2 — three network segments. Used by sections/03-platform.
% Monochrome line art: no fills, so it prints identically in colour and in
% greyscale. Keep it that way, and keep Fig. 3 matching it.
\begin{figure}[!t]
\centering
\resizebox{0.95\columnwidth}{!}{%
\begin{tikzpicture}[
  node distance=0.7cm and 1.5cm,
  box/.style={draw, rounded corners=3pt, align=center, font=\tiny,
              inner sep=4pt, minimum width=2.4cm},
  arr/.style={-{Stealth[length=3.5pt]}, semithick},
  lbl/.style={font=\tiny, fill=white, inner sep=1pt}
]
  \node[box] (pc) {\textbf{Development host}\\laboratory Wi-Fi};
  \node[box, below=0.7cm of pc] (head)
    {\textbf{Android head}\\wlan0: \texttt{172.16.0.x}\\eth0:\;\texttt{192.168.20.1}};
  \node[box, right=1.5cm of head] (ch)
    {\textbf{\ros{} chassis}\\WebSocket bridge :9090\\\texttt{192.168.20.22}};
  \node[box, below=0.7cm of ch] (ap)
    {\textbf{Robot access point}\\\texttt{10.42.0.1:9090}};
  \draw[arr] (pc)   -- node[lbl, left]{ADB} (head);
  \draw[arr] (head) -- node[lbl, above]{internal Ethernet} (ch);
  \draw[arr] (pc.east) .. controls +(1.0,-0.5) and +(-1.1,1.2) ..
             node[lbl, near start, right]{Wi-Fi} (ap.north);
  \draw[arr] (ap)   -- (ch);
\end{tikzpicture}%
}
\caption{Three network segments on a product marketed as a single device. Conflating the head's
DHCP address with the chassis address was the most frequent operator error we observed.}
\label{fig:network}
\end{figure}


Network mapping turned up three segments (Fig.~\ref{fig:network}), and that was already a
result in itself. Two channels are reachable from outside the vendor stack: a WebSocket bridge
on port~9090 and an MQTT broker on port~1883. Shell access to the chassis is closed. We observed
that repeated authentication attempts trigger an account lockout, and went no further down that
road, partly because it would break the non-destructive principle stated below and partly
because we had one unit and bricking it would have ended the study. Two vendor packages sit on
the head, a welcome application and a deployment tool, and both link against the same
chassis-communication library. We treat that library as the \emph{de facto} protocol
specification. It is the empirical basis for H1.
